% paper77-copilot-oracle-channel.tex — The Copilot Oracle Channel:
%   Federated Evidence from LLM and Solver Sources
%   Paper 77 of the Judgment Geometry series.
% Compile: pdflatex paper77-copilot-oracle-channel
\documentclass[11pt]{article}
\usepackage{lmodern}
% jugeo-common.tex — shared preamble for all Judgment Geometry papers
% Include via: % jugeo-common.tex — shared preamble for all Judgment Geometry papers
% Include via: % jugeo-common.tex — shared preamble for all Judgment Geometry papers
% Include via: \input{jugeo-common}

% ── Packages ────────────────────────────────────────────────────────────
\usepackage{amsmath,amssymb,amsthm,mathtools}
\usepackage{tikz-cd}
\usepackage{listings}
\usepackage{xcolor}
\usepackage{xspace}
\usepackage{booktabs}
\usepackage{multirow}
\usepackage{enumitem}
\usepackage[expansion=false]{microtype}
\usepackage{hyperref}
\usepackage{cleveref}
% \usepackage{thmtools}  % disabled: not available in this TeX installation
\usepackage{float}
\usepackage{caption}
\usepackage{subcaption}
\usepackage{tabularx}
\usepackage{stmaryrd}       % \llbracket, \rrbracket
\usepackage{bbm}            % \mathbbm{1}

% ── Page geometry ───────────────────────────────────────────────────────
\usepackage[margin=1in]{geometry}

% ── Theorem environments ────────────────────────────────────────────────
\theoremstyle{plain}
\newtheorem{theorem}{Theorem}[section]
\newtheorem{lemma}[theorem]{Lemma}
\newtheorem{proposition}[theorem]{Proposition}
\newtheorem{corollary}[theorem]{Corollary}

\theoremstyle{definition}
\newtheorem{definition}[theorem]{Definition}
\newtheorem{example}[theorem]{Example}
\newtheorem{construction}[theorem]{Construction}

\theoremstyle{remark}
\newtheorem{remark}[theorem]{Remark}
\newtheorem{notation}[theorem]{Notation}

% ── Python listings style ──────────────────────────────────────────────
\definecolor{jugeoBlue}{HTML}{2563EB}
\definecolor{jugeoGreen}{HTML}{059669}
\definecolor{jugeoRed}{HTML}{DC2626}
\definecolor{jugeoPurple}{HTML}{7C3AED}
\definecolor{jugeoGray}{HTML}{6B7280}
\definecolor{jugeoBg}{HTML}{F8FAFC}

\lstdefinestyle{jugeo-python}{
  language=Python,
  basicstyle=\ttfamily\small,
  keywordstyle=\color{jugeoBlue}\bfseries,
  stringstyle=\color{jugeoGreen},
  commentstyle=\color{jugeoGray}\itshape,
  numberstyle=\tiny\color{jugeoGray},
  backgroundcolor=\color{jugeoBg},
  numbers=left,
  numbersep=6pt,
  frame=single,
  framerule=0.4pt,
  rulecolor=\color{jugeoGray!40},
  breaklines=true,
  breakatwhitespace=true,
  showstringspaces=false,
  tabsize=4,
  xleftmargin=1.5em,
  framexleftmargin=1.5em,
  morekeywords={dataclass,frozen,field,Enum,IntEnum,auto,Any,
                Mapping,Sequence,Optional,match,case},
  literate={→}{{$\to$}}1 {←}{{$\leftarrow$}}1
           {≤}{{$\leq$}}1 {≥}{{$\geq$}}1 {≠}{{$\neq$}}1
           {∈}{{$\in$}}1 {∀}{{$\forall$}}1 {∃}{{$\exists$}}1
           {⊕}{{$\oplus$}}1 {⊖}{{$\ominus$}}1,
  escapeinside={(*@}{@*)},
}
\lstset{style=jugeo-python}

\lstdefinestyle{jugeo-output}{
  basicstyle=\ttfamily\small,
  backgroundcolor=\color{jugeoBg},
  frame=single,
  framerule=0.4pt,
  rulecolor=\color{jugeoGray!40},
  breaklines=true,
  showstringspaces=false,
  xleftmargin=1.5em,
  framexleftmargin=0.5em,
  numbers=none,
}

% ── JuGeo-specific macros ──────────────────────────────────────────────

% Judgment 8-tuple
\newcommand{\J}{\mathcal{J}}
\newcommand{\Jud}[8]{(#1,\,#2,\,#3,\,#4,\,#5,\,#6,\,#7,\,#8)}
\newcommand{\JudShort}[1]{\J_{#1}}

% Components
\newcommand{\coord}{\mathsf{c}}            % coordinate
\newcommand{\prop}{\varphi}                 % proposition
\newcommand{\carrier}{\mathcal{A}}          % carrier type
\newcommand{\evid}{\mathcal{E}}             % evidence bundle
\newcommand{\oblig}{\mathcal{O}}            % obligations
\newcommand{\obstr}{\mathcal{B}}            % obstructions
\newcommand{\trust}{\mathcal{T}}            % trust annotation
\newcommand{\prov}{\Pi}                     % provenance

% Trust algebra
\newcommand{\Talg}{\mathfrak{T}}            % trust ordered algebra
\newcommand{\Eadm}{\mathcal{E}_{\mathrm{adm}}}
\newcommand{\tleq}{\preceq}                 % trust ordering
\newcommand{\tjoin}{\oplus}                 % conservative join
\newcommand{\tatten}{\ominus}               % attenuation
\newcommand{\tpromote}{\uparrow_\pi}        % promotion
\newcommand{\tdemote}{\downarrow_\chi}       % demotion

% Trust tiers
\newcommand{\Tcontra}{\ensuremath{\mathsf{CONTRADICTED}}}
\newcommand{\Tunver}{\ensuremath{\mathsf{UNVERIFIED}}}
\newcommand{\Tcopilot}{\ensuremath{\mathsf{COPILOT}}}
\newcommand{\Toracle}{\ensuremath{\mathsf{ORACLE}}}
\newcommand{\Truntime}{\ensuremath{\mathsf{RUNTIME}}}
\newcommand{\Tsolver}{\ensuremath{\mathsf{SOLVER}}}
\newcommand{\Tproof}{\ensuremath{\mathsf{PROOF}}}

% Site and geometry
\newcommand{\Site}{\mathcal{S}}             % semantic site
\newcommand{\Cov}{\mathrm{Cov}}            % covering family
\newcommand{\Ob}{\mathrm{Ob}}              % objects
\newcommand{\Mor}{\mathrm{Mor}}            % morphisms
\newcommand{\res}{\mathrm{res}}            % restriction
\newcommand{\Cech}{\check{C}}              % Čech complex
\newcommand{\Hcoh}[1]{H^{#1}}             % cohomology group

% Descent
\newcommand{\descent}{\mathrm{desc}}
\newcommand{\glue}{\mathrm{glue}}
\newcommand{\overlap}[2]{#1 \cap #2}

% Semantic moves
\newcommand{\VERIFY}{\textsc{Verify}}
\newcommand{\CONSTRUCT}{\textsc{Construct}}
\newcommand{\NORMALIZE}{\textsc{Normalize}}
\newcommand{\GLUE}{\textsc{Glue}}
\newcommand{\RESTRICT}{\textsc{Restrict}}
\newcommand{\TRANSPORT}{\textsc{Transport}}
\newcommand{\REPAIR}{\textsc{Repair}}
\newcommand{\PROMOTE}{\textsc{Promote}}
\newcommand{\CHALLENGE}{\textsc{Challenge}}
\newcommand{\ESCALATE}{\textsc{Escalate}}

% Fragments
\newcommand{\QFLIA}{\texttt{QF\_LIA}}
\newcommand{\QFLRA}{\texttt{QF\_LRA}}
\newcommand{\QFBV}{\texttt{QF\_BV}}
\newcommand{\QFUF}{\texttt{QF\_UF}}

% Misc
\newcommand{\Python}{\textsc{Python}}
\newcommand{\JuGeo}{\textsc{JuGeo}}
\newcommand{\LEAN}{\textsc{Lean}\,4}
\newcommand{\Fstar}{F$^{\star}$}
\newcommand{\Coq}{\textsc{Coq}}
\newcommand{\Dafny}{\textsc{Dafny}}

% Common shortcuts
\newcommand{\ie}{i.e.\@\xspace}
\newcommand{\eg}{e.g.\@\xspace}
\newcommand{\cf}{cf.\@\xspace}
\newcommand{\wrt}{w.r.t.\@\xspace}

% Evaluation shortcuts
\newcommand{\numcases}{300}
\newcommand{\numspec}{100}
\newcommand{\numeq}{100}
\newcommand{\numbug}{100}
\newcommand{\medtime}{6.5\,ms}
\newcommand{\totaltime}{1.949\,s}


% ── Packages ────────────────────────────────────────────────────────────
\usepackage{amsmath,amssymb,amsthm,mathtools}
\usepackage{tikz-cd}
\usepackage{listings}
\usepackage{xcolor}
\usepackage{xspace}
\usepackage{booktabs}
\usepackage{multirow}
\usepackage{enumitem}
\usepackage[expansion=false]{microtype}
\usepackage{hyperref}
\usepackage{cleveref}
% \usepackage{thmtools}  % disabled: not available in this TeX installation
\usepackage{float}
\usepackage{caption}
\usepackage{subcaption}
\usepackage{tabularx}
\usepackage{stmaryrd}       % \llbracket, \rrbracket
\usepackage{bbm}            % \mathbbm{1}

% ── Page geometry ───────────────────────────────────────────────────────
\usepackage[margin=1in]{geometry}

% ── Theorem environments ────────────────────────────────────────────────
\theoremstyle{plain}
\newtheorem{theorem}{Theorem}[section]
\newtheorem{lemma}[theorem]{Lemma}
\newtheorem{proposition}[theorem]{Proposition}
\newtheorem{corollary}[theorem]{Corollary}

\theoremstyle{definition}
\newtheorem{definition}[theorem]{Definition}
\newtheorem{example}[theorem]{Example}
\newtheorem{construction}[theorem]{Construction}

\theoremstyle{remark}
\newtheorem{remark}[theorem]{Remark}
\newtheorem{notation}[theorem]{Notation}

% ── Python listings style ──────────────────────────────────────────────
\definecolor{jugeoBlue}{HTML}{2563EB}
\definecolor{jugeoGreen}{HTML}{059669}
\definecolor{jugeoRed}{HTML}{DC2626}
\definecolor{jugeoPurple}{HTML}{7C3AED}
\definecolor{jugeoGray}{HTML}{6B7280}
\definecolor{jugeoBg}{HTML}{F8FAFC}

\lstdefinestyle{jugeo-python}{
  language=Python,
  basicstyle=\ttfamily\small,
  keywordstyle=\color{jugeoBlue}\bfseries,
  stringstyle=\color{jugeoGreen},
  commentstyle=\color{jugeoGray}\itshape,
  numberstyle=\tiny\color{jugeoGray},
  backgroundcolor=\color{jugeoBg},
  numbers=left,
  numbersep=6pt,
  frame=single,
  framerule=0.4pt,
  rulecolor=\color{jugeoGray!40},
  breaklines=true,
  breakatwhitespace=true,
  showstringspaces=false,
  tabsize=4,
  xleftmargin=1.5em,
  framexleftmargin=1.5em,
  morekeywords={dataclass,frozen,field,Enum,IntEnum,auto,Any,
                Mapping,Sequence,Optional,match,case},
  literate={→}{{$\to$}}1 {←}{{$\leftarrow$}}1
           {≤}{{$\leq$}}1 {≥}{{$\geq$}}1 {≠}{{$\neq$}}1
           {∈}{{$\in$}}1 {∀}{{$\forall$}}1 {∃}{{$\exists$}}1
           {⊕}{{$\oplus$}}1 {⊖}{{$\ominus$}}1,
  escapeinside={(*@}{@*)},
}
\lstset{style=jugeo-python}

\lstdefinestyle{jugeo-output}{
  basicstyle=\ttfamily\small,
  backgroundcolor=\color{jugeoBg},
  frame=single,
  framerule=0.4pt,
  rulecolor=\color{jugeoGray!40},
  breaklines=true,
  showstringspaces=false,
  xleftmargin=1.5em,
  framexleftmargin=0.5em,
  numbers=none,
}

% ── JuGeo-specific macros ──────────────────────────────────────────────

% Judgment 8-tuple
\newcommand{\J}{\mathcal{J}}
\newcommand{\Jud}[8]{(#1,\,#2,\,#3,\,#4,\,#5,\,#6,\,#7,\,#8)}
\newcommand{\JudShort}[1]{\J_{#1}}

% Components
\newcommand{\coord}{\mathsf{c}}            % coordinate
\newcommand{\prop}{\varphi}                 % proposition
\newcommand{\carrier}{\mathcal{A}}          % carrier type
\newcommand{\evid}{\mathcal{E}}             % evidence bundle
\newcommand{\oblig}{\mathcal{O}}            % obligations
\newcommand{\obstr}{\mathcal{B}}            % obstructions
\newcommand{\trust}{\mathcal{T}}            % trust annotation
\newcommand{\prov}{\Pi}                     % provenance

% Trust algebra
\newcommand{\Talg}{\mathfrak{T}}            % trust ordered algebra
\newcommand{\Eadm}{\mathcal{E}_{\mathrm{adm}}}
\newcommand{\tleq}{\preceq}                 % trust ordering
\newcommand{\tjoin}{\oplus}                 % conservative join
\newcommand{\tatten}{\ominus}               % attenuation
\newcommand{\tpromote}{\uparrow_\pi}        % promotion
\newcommand{\tdemote}{\downarrow_\chi}       % demotion

% Trust tiers
\newcommand{\Tcontra}{\ensuremath{\mathsf{CONTRADICTED}}}
\newcommand{\Tunver}{\ensuremath{\mathsf{UNVERIFIED}}}
\newcommand{\Tcopilot}{\ensuremath{\mathsf{COPILOT}}}
\newcommand{\Toracle}{\ensuremath{\mathsf{ORACLE}}}
\newcommand{\Truntime}{\ensuremath{\mathsf{RUNTIME}}}
\newcommand{\Tsolver}{\ensuremath{\mathsf{SOLVER}}}
\newcommand{\Tproof}{\ensuremath{\mathsf{PROOF}}}

% Site and geometry
\newcommand{\Site}{\mathcal{S}}             % semantic site
\newcommand{\Cov}{\mathrm{Cov}}            % covering family
\newcommand{\Ob}{\mathrm{Ob}}              % objects
\newcommand{\Mor}{\mathrm{Mor}}            % morphisms
\newcommand{\res}{\mathrm{res}}            % restriction
\newcommand{\Cech}{\check{C}}              % Čech complex
\newcommand{\Hcoh}[1]{H^{#1}}             % cohomology group

% Descent
\newcommand{\descent}{\mathrm{desc}}
\newcommand{\glue}{\mathrm{glue}}
\newcommand{\overlap}[2]{#1 \cap #2}

% Semantic moves
\newcommand{\VERIFY}{\textsc{Verify}}
\newcommand{\CONSTRUCT}{\textsc{Construct}}
\newcommand{\NORMALIZE}{\textsc{Normalize}}
\newcommand{\GLUE}{\textsc{Glue}}
\newcommand{\RESTRICT}{\textsc{Restrict}}
\newcommand{\TRANSPORT}{\textsc{Transport}}
\newcommand{\REPAIR}{\textsc{Repair}}
\newcommand{\PROMOTE}{\textsc{Promote}}
\newcommand{\CHALLENGE}{\textsc{Challenge}}
\newcommand{\ESCALATE}{\textsc{Escalate}}

% Fragments
\newcommand{\QFLIA}{\texttt{QF\_LIA}}
\newcommand{\QFLRA}{\texttt{QF\_LRA}}
\newcommand{\QFBV}{\texttt{QF\_BV}}
\newcommand{\QFUF}{\texttt{QF\_UF}}

% Misc
\newcommand{\Python}{\textsc{Python}}
\newcommand{\JuGeo}{\textsc{JuGeo}}
\newcommand{\LEAN}{\textsc{Lean}\,4}
\newcommand{\Fstar}{F$^{\star}$}
\newcommand{\Coq}{\textsc{Coq}}
\newcommand{\Dafny}{\textsc{Dafny}}

% Common shortcuts
\newcommand{\ie}{i.e.\@\xspace}
\newcommand{\eg}{e.g.\@\xspace}
\newcommand{\cf}{cf.\@\xspace}
\newcommand{\wrt}{w.r.t.\@\xspace}

% Evaluation shortcuts
\newcommand{\numcases}{300}
\newcommand{\numspec}{100}
\newcommand{\numeq}{100}
\newcommand{\numbug}{100}
\newcommand{\medtime}{6.5\,ms}
\newcommand{\totaltime}{1.949\,s}


% ── Packages ────────────────────────────────────────────────────────────
\usepackage{amsmath,amssymb,amsthm,mathtools}
\usepackage{tikz-cd}
\usepackage{listings}
\usepackage{xcolor}
\usepackage{xspace}
\usepackage{booktabs}
\usepackage{multirow}
\usepackage{enumitem}
\usepackage[expansion=false]{microtype}
\usepackage{hyperref}
\usepackage{cleveref}
% \usepackage{thmtools}  % disabled: not available in this TeX installation
\usepackage{float}
\usepackage{caption}
\usepackage{subcaption}
\usepackage{tabularx}
\usepackage{stmaryrd}       % \llbracket, \rrbracket
\usepackage{bbm}            % \mathbbm{1}

% ── Page geometry ───────────────────────────────────────────────────────
\usepackage[margin=1in]{geometry}

% ── Theorem environments ────────────────────────────────────────────────
\theoremstyle{plain}
\newtheorem{theorem}{Theorem}[section]
\newtheorem{lemma}[theorem]{Lemma}
\newtheorem{proposition}[theorem]{Proposition}
\newtheorem{corollary}[theorem]{Corollary}

\theoremstyle{definition}
\newtheorem{definition}[theorem]{Definition}
\newtheorem{example}[theorem]{Example}
\newtheorem{construction}[theorem]{Construction}

\theoremstyle{remark}
\newtheorem{remark}[theorem]{Remark}
\newtheorem{notation}[theorem]{Notation}

% ── Python listings style ──────────────────────────────────────────────
\definecolor{jugeoBlue}{HTML}{2563EB}
\definecolor{jugeoGreen}{HTML}{059669}
\definecolor{jugeoRed}{HTML}{DC2626}
\definecolor{jugeoPurple}{HTML}{7C3AED}
\definecolor{jugeoGray}{HTML}{6B7280}
\definecolor{jugeoBg}{HTML}{F8FAFC}

\lstdefinestyle{jugeo-python}{
  language=Python,
  basicstyle=\ttfamily\small,
  keywordstyle=\color{jugeoBlue}\bfseries,
  stringstyle=\color{jugeoGreen},
  commentstyle=\color{jugeoGray}\itshape,
  numberstyle=\tiny\color{jugeoGray},
  backgroundcolor=\color{jugeoBg},
  numbers=left,
  numbersep=6pt,
  frame=single,
  framerule=0.4pt,
  rulecolor=\color{jugeoGray!40},
  breaklines=true,
  breakatwhitespace=true,
  showstringspaces=false,
  tabsize=4,
  xleftmargin=1.5em,
  framexleftmargin=1.5em,
  morekeywords={dataclass,frozen,field,Enum,IntEnum,auto,Any,
                Mapping,Sequence,Optional,match,case},
  literate={→}{{$\to$}}1 {←}{{$\leftarrow$}}1
           {≤}{{$\leq$}}1 {≥}{{$\geq$}}1 {≠}{{$\neq$}}1
           {∈}{{$\in$}}1 {∀}{{$\forall$}}1 {∃}{{$\exists$}}1
           {⊕}{{$\oplus$}}1 {⊖}{{$\ominus$}}1,
  escapeinside={(*@}{@*)},
}
\lstset{style=jugeo-python}

\lstdefinestyle{jugeo-output}{
  basicstyle=\ttfamily\small,
  backgroundcolor=\color{jugeoBg},
  frame=single,
  framerule=0.4pt,
  rulecolor=\color{jugeoGray!40},
  breaklines=true,
  showstringspaces=false,
  xleftmargin=1.5em,
  framexleftmargin=0.5em,
  numbers=none,
}

% ── JuGeo-specific macros ──────────────────────────────────────────────

% Judgment 8-tuple
\newcommand{\J}{\mathcal{J}}
\newcommand{\Jud}[8]{(#1,\,#2,\,#3,\,#4,\,#5,\,#6,\,#7,\,#8)}
\newcommand{\JudShort}[1]{\J_{#1}}

% Components
\newcommand{\coord}{\mathsf{c}}            % coordinate
\newcommand{\prop}{\varphi}                 % proposition
\newcommand{\carrier}{\mathcal{A}}          % carrier type
\newcommand{\evid}{\mathcal{E}}             % evidence bundle
\newcommand{\oblig}{\mathcal{O}}            % obligations
\newcommand{\obstr}{\mathcal{B}}            % obstructions
\newcommand{\trust}{\mathcal{T}}            % trust annotation
\newcommand{\prov}{\Pi}                     % provenance

% Trust algebra
\newcommand{\Talg}{\mathfrak{T}}            % trust ordered algebra
\newcommand{\Eadm}{\mathcal{E}_{\mathrm{adm}}}
\newcommand{\tleq}{\preceq}                 % trust ordering
\newcommand{\tjoin}{\oplus}                 % conservative join
\newcommand{\tatten}{\ominus}               % attenuation
\newcommand{\tpromote}{\uparrow_\pi}        % promotion
\newcommand{\tdemote}{\downarrow_\chi}       % demotion

% Trust tiers
\newcommand{\Tcontra}{\ensuremath{\mathsf{CONTRADICTED}}}
\newcommand{\Tunver}{\ensuremath{\mathsf{UNVERIFIED}}}
\newcommand{\Tcopilot}{\ensuremath{\mathsf{COPILOT}}}
\newcommand{\Toracle}{\ensuremath{\mathsf{ORACLE}}}
\newcommand{\Truntime}{\ensuremath{\mathsf{RUNTIME}}}
\newcommand{\Tsolver}{\ensuremath{\mathsf{SOLVER}}}
\newcommand{\Tproof}{\ensuremath{\mathsf{PROOF}}}

% Site and geometry
\newcommand{\Site}{\mathcal{S}}             % semantic site
\newcommand{\Cov}{\mathrm{Cov}}            % covering family
\newcommand{\Ob}{\mathrm{Ob}}              % objects
\newcommand{\Mor}{\mathrm{Mor}}            % morphisms
\newcommand{\res}{\mathrm{res}}            % restriction
\newcommand{\Cech}{\check{C}}              % Čech complex
\newcommand{\Hcoh}[1]{H^{#1}}             % cohomology group

% Descent
\newcommand{\descent}{\mathrm{desc}}
\newcommand{\glue}{\mathrm{glue}}
\newcommand{\overlap}[2]{#1 \cap #2}

% Semantic moves
\newcommand{\VERIFY}{\textsc{Verify}}
\newcommand{\CONSTRUCT}{\textsc{Construct}}
\newcommand{\NORMALIZE}{\textsc{Normalize}}
\newcommand{\GLUE}{\textsc{Glue}}
\newcommand{\RESTRICT}{\textsc{Restrict}}
\newcommand{\TRANSPORT}{\textsc{Transport}}
\newcommand{\REPAIR}{\textsc{Repair}}
\newcommand{\PROMOTE}{\textsc{Promote}}
\newcommand{\CHALLENGE}{\textsc{Challenge}}
\newcommand{\ESCALATE}{\textsc{Escalate}}

% Fragments
\newcommand{\QFLIA}{\texttt{QF\_LIA}}
\newcommand{\QFLRA}{\texttt{QF\_LRA}}
\newcommand{\QFBV}{\texttt{QF\_BV}}
\newcommand{\QFUF}{\texttt{QF\_UF}}

% Misc
\newcommand{\Python}{\textsc{Python}}
\newcommand{\JuGeo}{\textsc{JuGeo}}
\newcommand{\LEAN}{\textsc{Lean}\,4}
\newcommand{\Fstar}{F$^{\star}$}
\newcommand{\Coq}{\textsc{Coq}}
\newcommand{\Dafny}{\textsc{Dafny}}

% Common shortcuts
\newcommand{\ie}{i.e.\@\xspace}
\newcommand{\eg}{e.g.\@\xspace}
\newcommand{\cf}{cf.\@\xspace}
\newcommand{\wrt}{w.r.t.\@\xspace}

% Evaluation shortcuts
\newcommand{\numcases}{300}
\newcommand{\numspec}{100}
\newcommand{\numeq}{100}
\newcommand{\numbug}{100}
\newcommand{\medtime}{6.5\,ms}
\newcommand{\totaltime}{1.949\,s}

% No external data file for this paper.

% ── Additional packages ────────────────────────────────────────────
\usepackage{array}
\usepackage{tikz}
\usetikzlibrary{arrows.meta,positioning,calc,fit,backgrounds}

% ── Paper-specific macros ──────────────────────────────────────────
\newcommand{\CopOracle}{\mathsf{CopilotOracleChannel}}
\newcommand{\OracleComp}{\mathrm{compose}}
\newcommand{\OracleMerge}{\mathrm{merge}}
\newcommand{\OracleValid}{\mathrm{validate}}
\newcommand{\OracleEscal}{\mathrm{escalate}}
\newcommand{\FedResult}{\mathsf{FederationResult}}
\newcommand{\CopChan}{\mathsf{CopilotChannel}}
\newcommand{\SolChan}{\mathsf{SolverChannel}}
\newcommand{\ChanFed}{\mathsf{ChannelFederation}}
\newcommand{\ChanRouter}{\mathsf{ChannelRouter}}
\newcommand{\EvReq}{\mathsf{EvidenceRequest}}
\newcommand{\EvResp}{\mathsf{EvidenceResponse}}
\newcommand{\EvRec}{\mathsf{EvidenceRecord}}
\newcommand{\SupRoute}{\mathsf{SupportRoute}}
\newcommand{\TrustCeil}{\mathrm{ceiling}}
\newcommand{\NoEscal}{\mathrm{noEscalation}}

\title{\textbf{The Copilot Oracle Channel:\\
Federated Evidence from LLM and Solver Sources}}

\author{Halley Young}

\date{Paper~77 --- Judgment Geometry Series}

\begin{document}
\maketitle

% ─────────────────────────────────────────────────────────────────────
\begin{abstract}
Modern program verification benefits from combining evidence from
multiple sources: SMT solvers provide sound proofs, runtime monitors
capture execution witnesses, and large language models propose
plausible conjectures.  Combining these heterogeneous sources into
a coherent evidence picture while maintaining formal trust guarantees
is a central challenge.  We introduce the \emph{Copilot Oracle
Channel} ($\CopOracle$), a federated evidence architecture within
\JuGeo{} that composes evidence from the $\CopChan$ and the
$\SolChan$ through the $\ChanFed$ mechanism.

The federation operates under a strict no-escalation invariant:
the trust level of composed evidence never exceeds the minimum
trust ceiling of the contributing channels.  We formalize
federation as a meet operation on the trust lattice, prove that
composition preserves the no-escalation property, and establish
bounds on the effectiveness of federated evidence relative to
single-channel evidence.

We reference \pSevenSevenWorkflows{} agentic workflows spanning
\pSevenSevenSteps{} verification steps and
\pSevenSevenToolCalls{} tool invocations.  The federated approach
achieves a step-level verification rate of
\pSevenSevenStepVerifRate{} with an obstruction detection rate of
\pSevenSevenDiagnosisRate.  All formal properties are proved in
full and formalized in \LEAN.
\end{abstract}

% ─────────────────────────────────────────────────────────────────────
\section{Introduction}
\label{sec:intro}

Formal verification at scale requires drawing on multiple evidence
sources.  A solver may efficiently discharge arithmetic obligations
but struggle with heap reasoning; a runtime monitor may capture
witnesses that are difficult to construct symbolically; an LLM
may suggest high-level strategies that neither solvers nor monitors
can discover.  The challenge is to combine these diverse sources
while maintaining the trust guarantees that each source provides.

In \JuGeo{}, the evidence channel architecture provides a structured
framework for multi-source evidence.  Each channel operates within
its trust interval, and the $\ChanFed$ mechanism composes evidence
from multiple channels.  The $\CopOracle$ extends this architecture
to the specific case of Copilot-solver federation, where the
complementary strengths of LLM inference and SMT solving are
combined.

\paragraph{Motivation.}
The Copilot channel excels at generating plausible hypotheses
(invariants, strategies, rewrite suggestions) but cannot verify
them.  The solver channel excels at verification but may struggle
with hypothesis generation.  By federating the two, we obtain a
system that generates and verifies hypotheses within a single
trust-calibrated pipeline.

\paragraph{Contributions.}
\begin{itemize}[leftmargin=2em]
\item We define the $\CopOracle$ as a federated channel that
  composes Copilot and solver evidence (\S\ref{sec:oracle}).

\item We formalize the no-escalation invariant and prove it is
  preserved under arbitrary federation sequences (\S\ref{sec:formal}).

\item We introduce support routes as the mechanism for linking
  Copilot proposals to solver validations (\S\ref{sec:routes}).

\item We prove five formal properties: no-escalation, federation
  commutativity, federation associativity, support route soundness,
  and router correctness (\S\ref{sec:formal}).

\item We describe the integration with the $\ChanRouter$ for
  fallback-aware routing (\S\ref{sec:routing}).
\end{itemize}

\paragraph{Outline.}
Section~\ref{sec:background} reviews evidence channels.
Section~\ref{sec:oracle} defines the oracle channel.
Section~\ref{sec:routes} formalizes support routes.
Section~\ref{sec:routing} describes routing.
Section~\ref{sec:formal} proves the main theorems.
Section~\ref{sec:discussion} discusses implications.
Section~\ref{sec:related} surveys related work.
Section~\ref{sec:conclusion} concludes.

% ─────────────────────────────────────────────────────────────────────
\section{Background: Evidence Federation}
\label{sec:background}

\subsection{Channel Trust Model}
\label{sec:background:trust}

Each evidence channel $C$ operates within a trust interval
$I(C) = [\tau_{\min}(C), \tau_{\max}(C)]$.  The key intervals are:
\begin{align*}
  I(\mathsf{COPILOT}) &= [\Tcopilot, \Tcopilot], \\
  I(\mathsf{SOLVER}) &= [\Tsolver, \Tproof], \\
  I(\mathsf{RUNTIME}) &= [\Truntime, \Tproof], \\
  I(\mathsf{ORACLE}) &= [\Toracle, \Tproof].
\end{align*}

\begin{definition}[Trust composition]
\label{def:trust-comp}
The trust of composed evidence from channels $C_1, \ldots, C_n$
is the meet of the individual trusts:
\[
  \tau^* = \tau_1 \tjoin \tau_2 \tjoin \cdots \tjoin \tau_n
  = \min(\tau_1, \ldots, \tau_n).
\]
\end{definition}

\begin{remark}[Conservative composition]
\label{rem:conservative}
The meet-based composition is inherently conservative: the
composed trust is never higher than any individual contribution.
This prevents a low-trust channel from ``lifting'' its evidence
by associating with a high-trust channel.
\end{remark}

\subsection{The Federation Mechanism}
\label{sec:background:federation}

The $\ChanFed$ collects evidence from multiple channels, merges
it by evidence kind, and verifies the no-escalation property.

\begin{definition}[Federation]
\label{def:federation}
A \emph{federation} $\mathcal{F} = (C_1, \ldots, C_n)$ is an
ordered tuple of channels.  The federation produces a
$\FedResult$ containing:
\begin{itemize}[leftmargin=2em]
\item Merged evidence items organized by evidence kind.
\item A composed trust level $\tau^*$.
\item A verification that $\NoEscal$ holds.
\end{itemize}
\end{definition}

The following listing shows the core federation of Copilot and
solver channels.

\begin{lstlisting}[style=jugeo-python, caption={Federation of CopilotChannel and SolverChannel.}, label=lst:federation]
from jugeo.evidence.channels import (
    CopilotChannel,
    SolverChannel,
    ChannelFederation,
    ChannelPool,
    ChannelConfiguration,
    EvidenceChannel,
    EvidenceRequest,
    EvidenceResponse,
    EvidenceKind,
)

# Configure and instantiate both channels
cop_config = ChannelConfiguration.default_for(EvidenceChannel.COPILOT)
sol_config = ChannelConfiguration.default_for(EvidenceChannel.SOLVER)

copilot = CopilotChannel(config=cop_config, model_name="copilot-default")
solver = SolverChannel(config=sol_config)

# Register channels in the pool
pool = ChannelPool()
pool.register_channel(copilot)
pool.register_channel(solver)

# Create the federation
federation = ChannelFederation()

# Build a federated verification request
request = EvidenceRequest(
    request_id="oracle-fed-001",
    coordinate="binary_search.invariant",
    proposition="lo <= mid <= hi and arr[mid] == target",
    required_kind="PROPOSAL",
    proposition_kind="invariant",
    preferred_channel=EvidenceChannel.COPILOT,
    fallback_channels=(EvidenceChannel.SOLVER,),
    deadline_ms=15000.0,
    budget=3.0,
    metadata={"federation": True, "channels": ["copilot", "solver"]},
)

# Step 1: Get proposal from Copilot
cop_response = copilot.handle_request(request)
assert cop_response.trust_level == "proposal"

# Step 2: Validate with solver
sol_request = EvidenceRequest(
    request_id="oracle-fed-001-validate",
    coordinate=request.coordinate,
    proposition=request.proposition,
    required_kind="SOLVER",
    proposition_kind="validation",
    preferred_channel=EvidenceChannel.SOLVER,
    fallback_channels=(),
    deadline_ms=10000.0,
    budget=2.0,
    metadata={"copilot_evidence": cop_response.evidence},
)
sol_response = solver.handle_request(sol_request)

# Step 3: Federate the evidence
all_evidence = federation.collect_all()
merged = federation.merge_by_kind()

# Verify no trust escalation
federation.verify_no_escalation()

# Compose the final trust level
composed_trust = federation.compose_trust()
print(f"Copilot trust: {cop_response.trust_level}")
print(f"Solver trust: {sol_response.trust_level}")
print(f"Composed trust: {composed_trust}")

# Federation log for audit
fed_log = federation.federation_log()
print(f"Federation log entries: {len(fed_log)}")
federation.clear_log()
\end{lstlisting}

% ─────────────────────────────────────────────────────────────────────
\section{The Copilot Oracle Channel}
\label{sec:oracle}

\subsection{Design Rationale}
\label{sec:oracle:rationale}

The $\CopOracle$ represents a pattern of usage rather than a
single class: it is the composition of the $\CopChan$ and
$\SolChan$ through the $\ChanFed$, where the Copilot
generates hypotheses and the solver validates them.

\begin{definition}[Copilot oracle pattern]
\label{def:cop-oracle}
The \emph{Copilot oracle pattern} is the evidence production
strategy $\Pi_{\CopOracle}$ defined as:
\[
  \Pi_{\CopOracle}(q) = \OracleValid(\SolChan,\; \CopChan.\texttt{handle}(q))
\]
where $\OracleValid$ checks the Copilot's proposal against the
solver.  The result trust is:
\[
  \tau(\Pi_{\CopOracle}(q)) = \begin{cases}
    \Tsolver & \text{if solver confirms}, \\
    \Tcopilot & \text{if solver times out}, \\
    \Tcontra & \text{if solver refutes}.
  \end{cases}
\]
\end{definition}

\begin{remark}[Trust promotion through validation]
\label{rem:trust-promotion}
The oracle pattern is the primary mechanism for promoting
Copilot evidence above $\Tcopilot$.  When the solver confirms
a Copilot proposal, the \emph{combined} evidence (proposal +
confirmation) may receive trust $\Tsolver$, because the solver's
confirmation is the operative evidence.  The Copilot's role is
heuristic guidance, not evidentiary contribution to the final
trust level.
\end{remark}

\subsection{Oracle Composition}
\label{sec:oracle:composition}

\begin{definition}[Oracle composition operator]
\label{def:oracle-comp}
The oracle composition operator $\OracleComp$ takes two channel
responses $s_1, s_2$ and produces a merged response:
\[
  \OracleComp(s_1, s_2) = (w^*, \tau^*, \rho^*, \ell^*)
\]
where:
\begin{itemize}[leftmargin=2em]
\item $w^* = w_1 \cup w_2$ (union of witnesses),
\item $\tau^* = \tau_1 \tjoin \tau_2$ (trust meet),
\item $\rho^* = \rho_1 \cap \rho_2$ (intersection of residuals),
\item $\ell^* = \max(\ell_1, \ell_2)$ (maximum latency).
\end{itemize}
\end{definition}

\begin{lemma}[Composition trust bound]
\label{lem:comp-trust}
$\tau(\OracleComp(s_1, s_2)) \tleq \min(\tau(s_1), \tau(s_2))$.
\end{lemma}

\begin{proof}
By definition, $\tau^* = \tau_1 \tjoin \tau_2 = \min(\tau_1, \tau_2)$.
\end{proof}

\subsection{The No-Escalation Invariant}
\label{sec:oracle:noescal}

\begin{definition}[No-escalation invariant]
\label{def:no-escal}
A federation $\mathcal{F} = (C_1, \ldots, C_n)$ satisfies
\emph{no-escalation} if the composed trust $\tau^*$ satisfies:
\[
  \NoEscal(\mathcal{F}) \iff
  \tau^* \tleq \min_{i=1}^{n} \TrustCeil(C_i).
\]
\end{definition}

\begin{remark}[Why no-escalation matters]
\label{rem:why-noescal}
Without the no-escalation invariant, a malicious or buggy channel
could exploit federation to inflate the trust of its evidence.
For example, if the Copilot channel produced evidence labeled
$\Tsolver$ and the federation did not check ceilings, the
composed result could inherit this inflated trust.  The
no-escalation check prevents this by verifying that the composed
trust respects all ceilings.
\end{remark}

% ─────────────────────────────────────────────────────────────────────
\section{Support Routes}
\label{sec:routes}

Support routes formalize the relationship between a Copilot
proposal and its solver validation.

\subsection{Route Structure}
\label{sec:routes:structure}

\begin{definition}[Support route]
\label{def:support-route}
A \emph{support route} $\sigma$ is a tuple
$(k, t, w_s, w_f, p, r, q, d)$ where:
\begin{itemize}[leftmargin=2em]
\item $k$ is the channel kind (e.g., $\mathsf{PROPOSAL}$).
\item $t$ is the target evidence form.
\item $w_s$ is the success witness schema.
\item $w_f$ is the failure witness schema.
\item $p$ is the invalidation policy.
\item $r$ is the support region (coordinate prefix).
\item $q$ is the admissible query family.
\item $d$ is the human-readable reason string.
\end{itemize}
\end{definition}

\begin{definition}[Route validity]
\label{def:route-valid}
A support route $\sigma$ for evidence record $e$ is \emph{valid}
if:
\begin{enumerate}
\item $e.\mathrm{channel\_kind}$ matches $\sigma.k$.
\item $e.\mathrm{coordinate}$ is within $\sigma.r$.
\item The evidence form matches $\sigma.t$.
\end{enumerate}
\end{definition}

The following listing shows how evidence records are constructed
with support routes in the oracle pattern.

\begin{lstlisting}[style=jugeo-python, caption={Evidence records with support routes for oracle federation.}, label=lst:routes]
from jugeo.evidence.channels import (
    EvidenceRecord,
    SupportRoute,
    EvidenceChannel,
    EvidenceKind,
    ClaimPolarity,
    CopilotChannel,
    SolverChannel,
    ChannelConfiguration,
    EvidenceRequest,
)

# Set up the oracle pattern channels
copilot = CopilotChannel(
    config=ChannelConfiguration.default_for(EvidenceChannel.COPILOT)
)
solver = SolverChannel(
    config=ChannelConfiguration.default_for(EvidenceChannel.SOLVER)
)

# Define the support route: Copilot proposal -> Solver validation
oracle_route = SupportRoute(
    channel_kind=EvidenceKind.PROPOSAL,
    target_evidence_form="solver_certificate",
    success_witness_schema="z3_proof",
    failure_witness_schema="z3_counterexample",
    invalidation_policy="on_counterexample",
    support_region="binary_search",
    admissible_query_family="invariant_validation",
    reason="Copilot invariant proposal validated by Z3 solver",
)

# Create the evidence record for the Copilot proposal
copilot_record = EvidenceRecord(
    channel=EvidenceChannel.COPILOT,
    claim="lo <= mid <= hi",
    record_id="oracle-record-001",
    evidence={"type": "invariant", "formula": "lo <= mid <= hi"},
    trust_level="proposal",
    coordinate="binary_search.loop.invariant",
    payload={"model": "copilot-default", "confidence": 0.85},
    obligations=("validate_with_solver",),
    provenance=("copilot-query-001",),
)

# Attach the support route
copilot_record = copilot_record.with_support_route(oracle_route)

# Inspect the record
print(f"Record ID: {copilot_record.record_id}")
print(f"Trust level: {copilot_record.trust_level}")
print(f"Canonical key: {copilot_record.canonical_key()}")
print(f"Support regions: {copilot_record.support_regions()}")
print(f"Obligations: {copilot_record.obligations}")

# Build the solver validation request
validation_request = EvidenceRequest(
    request_id="oracle-validate-001",
    coordinate="binary_search.loop.invariant",
    proposition="lo <= mid <= hi",
    required_kind="SOLVER",
    proposition_kind="validation",
    preferred_channel=EvidenceChannel.SOLVER,
    fallback_channels=(),
    deadline_ms=10000.0,
    budget=2.0,
    metadata={
        "copilot_record_id": copilot_record.record_id,
        "support_route": oracle_route.canonical_key(),
    },
)

# Get solver validation
solver_response = solver.handle_request(validation_request)

# Create solver evidence record
solver_record = EvidenceRecord(
    channel=EvidenceChannel.SOLVER,
    claim="lo <= mid <= hi is valid",
    record_id="oracle-solver-001",
    evidence=solver_response.evidence,
    trust_level=solver_response.trust_level,
    coordinate="binary_search.loop.invariant",
    payload={"validated_record": copilot_record.record_id},
    obligations=(),
    provenance=("solver-session-001",),
)

# Check canonical keys for linking
cop_key = copilot_record.canonical_key()
sol_key = solver_record.canonical_key()
route_key = oracle_route.canonical_key()
print(f"Copilot key: {cop_key}")
print(f"Solver key: {sol_key}")
print(f"Route key: {route_key}")

# Serialize the route
route_dict = oracle_route.to_dict()
print(f"Route serialized: {list(route_dict.keys())}")
\end{lstlisting}

% ─────────────────────────────────────────────────────────────────────
\section{Routing and Fallback}
\label{sec:routing}

The $\ChanRouter$ manages the routing of evidence requests
across federated channels.

\subsection{Routing Strategy}
\label{sec:routing:strategy}

\begin{definition}[Oracle routing strategy]
\label{def:oracle-routing}
The \emph{oracle routing strategy} proceeds as follows:
\begin{enumerate}
\item Route the request to the $\CopChan$ for proposal generation.
\item If the $\CopChan$ produces a non-empty proposal, route a
  validation request to the $\SolChan$.
\item If the $\SolChan$ confirms, compose the evidence.
\item If the $\SolChan$ refutes, record the obstruction and
  optionally retry with a refined proposal.
\item If the $\CopChan$ fails or times out, fall back to the
  $\SolChan$ for direct evidence production.
\end{enumerate}
\end{definition}

\begin{definition}[Fallback chain]
\label{def:fallback}
A \emph{fallback chain} is an ordered list of channels
$[C_1, C_2, \ldots, C_k]$ where each channel serves as a
backup for the preceding one.  The router tries channels in
order until one succeeds.
\end{definition}

The following listing demonstrates the full federated verification
pipeline with routing and fallback.

\begin{lstlisting}[style=jugeo-python, caption={Full oracle verification pipeline with routing.}, label=lst:pipeline]
from jugeo.evidence.channels import (
    CopilotChannel,
    SolverChannel,
    RuntimeChannel,
    ChannelRouter,
    ChannelPool,
    ChannelMonitor,
    ChannelFederation,
    ChannelConfiguration,
    EvidenceChannel,
    EvidenceRequest,
    EvidenceResponse,
)
from jugeo.interfaces.api import (
    JuGeoAPI,
    APIRequest,
    APIResponse,
    OperationKind,
    RequestStatus,
    CopilotAPIBridge,
    APIEventLog,
)

# Build the full channel infrastructure
pool = ChannelPool()
copilot = CopilotChannel(
    config=ChannelConfiguration.default_for(EvidenceChannel.COPILOT)
)
solver = SolverChannel(
    config=ChannelConfiguration.default_for(EvidenceChannel.SOLVER)
)
pool.register_channel(copilot)
pool.register_channel(solver)

# Set up routing, monitoring, and federation
router = ChannelRouter()
monitor = ChannelMonitor()
federation = ChannelFederation()

# Set up API with Copilot bridge
event_log = APIEventLog()
bridge = CopilotAPIBridge(
    channel=copilot,
    event_log=event_log,
    require_corroboration=True,
)
api = JuGeoAPI(event_log=event_log, copilot_bridge=bridge)

# Oracle verification request
request = APIRequest(
    request_id="oracle-verify-001",
    operation=OperationKind.VERIFY,
    coordinate="quicksort.partition.postcondition",
    proposition="all(a[i] <= pivot for i in range(lo, j))",
    budget=25000,
    deadline=30.0,
    metadata={"strategy": "oracle-federated"},
)

# Execute federated verification
response = api.verify(request)

if response.is_successful():
    print(f"Verified at trust: {response.trust_level}")
    print(f"Copilot contributed: {response.copilot_contributed}")
elif response.has_residuals():
    for r in response.residuals:
        print(f"Residual: {r}")
else:
    print(f"Failed: {response.status}")

# Check federation integrity
federation.verify_no_escalation()
composed = federation.compose_trust()
print(f"Federated trust: {composed}")

# Router finds best channel for remaining work
best = router.find_best_channel()
all_capable = router.find_all_capable()
print(f"Best channel: {best}")
print(f"All capable: {all_capable}")

# Monitor performance
print(f"Success rate: {monitor.success_rate()}")
print(f"Throughput: {monitor.throughput()}")

# Copilot bridge statistics
stats = bridge.stats()
print(f"Bridge stats: {stats}")

# Event log
events = event_log.export()
print(f"Events: {len(event_log)}")
\end{lstlisting}

\subsection{Router Correctness}
\label{sec:routing:correct}

\begin{definition}[Router correctness]
\label{def:router-correct}
The router $R$ is \emph{correct} if for every request $q$:
\begin{enumerate}
\item If $R(q) = C$, then $C$'s jurisdiction admits $q$.
\item If $R(q) = \bot$ (no channel selected), then no available
  channel admits $q$.
\item The selected channel has the highest priority among
  admissible channels.
\end{enumerate}
\end{definition}

% ─────────────────────────────────────────────────────────────────────
\section{Formal Properties}
\label{sec:formal}

We prove five formal properties.  All proofs are formalized in
\LEAN{} in \texttt{Paper77\_CopilotOracle.lean}.

\subsection{No-Escalation Theorem}
\label{sec:formal:noescal}

\begin{theorem}[No-escalation under federation]
\label{thm:no-escal}
For any federation $\mathcal{F} = (C_1, \ldots, C_n)$ with
trust ceilings $\TrustCeil(C_1), \ldots, \TrustCeil(C_n)$,
the composed trust satisfies:
\[
  \tau^* \tleq \min_{i=1}^{n} \TrustCeil(C_i).
\]
In particular, if $\mathcal{F}$ contains the $\CopChan$:
$\tau^* \tleq \Tcopilot$.
\end{theorem}

\begin{proof}
The composed trust is $\tau^* = \tau_1 \tjoin \cdots \tjoin \tau_n$.
Each $\tau_i \tleq \TrustCeil(C_i)$ by the trust ceiling of
channel $C_i$.  Since $\tjoin$ computes the minimum:
\[
  \tau^* = \min(\tau_1, \ldots, \tau_n)
  \leq \min(\TrustCeil(C_1), \ldots, \TrustCeil(C_n)).
\]
If any $C_i$ is the $\CopChan$, then
$\TrustCeil(C_i) = \Tcopilot$, so
$\min_j \TrustCeil(C_j) \leq \Tcopilot$, and thus
$\tau^* \tleq \Tcopilot$.

The $\ChanFed.\texttt{verify\_no\_escalation}()$ method
checks this condition post-hoc and raises an error if violated,
providing a runtime safety net in addition to the compile-time
guarantee.
\end{proof}

\begin{corollary}[Copilot cannot escalate through federation]
\label{cor:cop-noescal}
No federation involving the $\CopChan$ can produce evidence
with trust above $\Tcopilot$, regardless of other participants.
\end{corollary}

\begin{proof}
Immediate from Theorem~\ref{thm:no-escal} with the $\CopChan$
as one of the federation members.
\end{proof}

\subsection{Federation Commutativity}
\label{sec:formal:comm}

\begin{theorem}[Federation is commutative]
\label{thm:fed-comm}
For any two channels $C_1, C_2$ with responses $s_1, s_2$:
$\OracleComp(s_1, s_2) = \OracleComp(s_2, s_1)$.
\end{theorem}

\begin{proof}
By Definition~\ref{def:oracle-comp}:
\begin{itemize}
\item $w^* = w_1 \cup w_2 = w_2 \cup w_1$ (set union is commutative).
\item $\tau^* = \tau_1 \tjoin \tau_2 = \tau_2 \tjoin \tau_1$
  ($\tjoin = \min$ is commutative).
\item $\rho^* = \rho_1 \cap \rho_2 = \rho_2 \cap \rho_1$
  (intersection is commutative).
\item $\ell^* = \max(\ell_1, \ell_2) = \max(\ell_2, \ell_1)$
  ($\max$ is commutative).
\end{itemize}
All components are equal, so the compositions are equal.
\end{proof}

\subsection{Federation Associativity}
\label{sec:formal:assoc}

\begin{theorem}[Federation is associative]
\label{thm:fed-assoc}
For responses $s_1, s_2, s_3$:
$\OracleComp(\OracleComp(s_1, s_2), s_3)
= \OracleComp(s_1, \OracleComp(s_2, s_3))$.
\end{theorem}

\begin{proof}
Each component operation ($\cup$, $\min$, $\cap$, $\max$)
is associative.  Therefore the componentwise composition is
associative.

Explicitly for trust:
$(\tau_1 \tjoin \tau_2) \tjoin \tau_3
= \min(\min(\tau_1, \tau_2), \tau_3)
= \min(\tau_1, \min(\tau_2, \tau_3))
= \tau_1 \tjoin (\tau_2 \tjoin \tau_3)$.

Similar arguments apply to each component.
\end{proof}

\begin{corollary}[Federation forms a commutative monoid]
\label{cor:fed-monoid}
The set of federation results under $\OracleComp$ with the identity
element $(\emptyset, \Tproof, \mathcal{U}, 0)$ forms a commutative
monoid.
\end{corollary}

\begin{proof}
Commutativity (Theorem~\ref{thm:fed-comm}), associativity
(Theorem~\ref{thm:fed-assoc}), and identity
($\OracleComp(s, (\emptyset, \Tproof, \mathcal{U}, 0)) = s$
since $\cup \emptyset = \mathrm{id}$, $\min(\tau, \Tproof) = \tau$,
$\cap \mathcal{U} = \mathrm{id}$, $\max(\ell, 0) = \ell$).
\end{proof}

\subsection{Support Route Soundness}
\label{sec:formal:routes}

\begin{theorem}[Support route soundness]
\label{thm:route-sound}
If an evidence record $e$ has a valid support route $\sigma$
(Definition~\ref{def:route-valid}) and the solver validates
the claim at the end of the route, then the composed evidence
achieves trust at least $\Tsolver$.
\end{theorem}

\begin{proof}
A valid support route links a Copilot proposal (trust $\Tcopilot$)
to a solver validation (trust $\Tsolver$).  When the solver
confirms the claim, the \emph{solver evidence} is the operative
component, not the Copilot proposal.  The composed trust is
determined by the solver's trust floor:
\[
  \tau^* = \Tsolver
\]
because the solver's independent confirmation does not depend on
the Copilot's trust level.

Note that this is \emph{not} a trust escalation of the Copilot
evidence; rather, the solver produces \emph{new} evidence (a proof
certificate) that independently supports the claim.  The Copilot's
contribution is the hypothesis, not the evidence.
\end{proof}

\subsection{Router Correctness}
\label{sec:formal:router}

\begin{theorem}[Router selects admissible channels]
\label{thm:router}
For every request $q$, if the router selects channel $C$, then
$C$'s jurisdiction admits $q$.
\end{theorem}

\begin{proof}
The $\ChanRouter.\texttt{route}()$ method first calls
$\texttt{check\_jurisdiction}(q)$ for each candidate channel.
Only channels that pass the jurisdiction check are added to the
candidate set.  The router then selects the highest-priority
candidate.  Since only admissible channels are candidates, the
selected channel's jurisdiction admits $q$.
\end{proof}

% ─────────────────────────────────────────────────────────────────────
\section{Discussion}
\label{sec:discussion}

\paragraph{Copilot as hypothesis generator.}
The oracle pattern assigns a clear role to the Copilot channel:
hypothesis generation.  The Copilot excels at suggesting plausible
candidates (invariants, postconditions, lemma applications) based
on pattern matching.  The solver then serves as the arbiter of
correctness.  This division of labor leverages the complementary
strengths of both systems.

\paragraph{Latency implications.}
The oracle pattern introduces a sequential dependency: the solver
cannot begin validation until the Copilot produces a proposal.
The total latency is approximately $\ell_{\mathrm{Copilot}} +
\ell_{\mathrm{Solver}}$.  In practice, the Copilot's latency
(~200--500ms for a single query) is comparable to the solver's
(~100--5000ms depending on the formula), so the overhead is
moderate.

\paragraph{Failure modes.}
The oracle pattern can fail in several ways: the Copilot may
produce unhelpful proposals, the solver may time out on valid
proposals, or the solver may incorrectly refute a correct proposal
due to theory incompleteness.  The fallback mechanism
(Definition~\ref{def:fallback}) provides resilience against
these failures.

\paragraph{Multi-solver federation.}
The oracle pattern generalizes to federations involving multiple
solvers: a Copilot proposal can be validated by Z3, CVC5, or a
custom decision procedure.  The federation mechanism handles this
seamlessly since all solvers share the same trust ceiling.

% ─────────────────────────────────────────────────────────────────────
\section{Related Work}
\label{sec:related}

\paragraph{Oracle-guided verification.}
The use of oracles in verification has a long
history~\cite{hoare1969axiomatic,floyd1967assigning}.
In model checking, oracles guide abstraction
refinement~\cite{clarke2003counterexample}.  Our oracle pattern
adapts this concept to the LLM setting.

\paragraph{Multi-engine verification.}
Combining multiple verification engines has been explored in
portfolio solvers~\cite{xu2008satzilla} and parallel
verification~\cite{wintersteiger2009portfolio}.  Our
contribution is the trust-calibrated federation mechanism.

\paragraph{LLM-solver integration.}
Recent work has explored combining LLMs with SMT
solvers~\cite{first2023baldur,yang2024leandojo}.  Our framework
provides the formal foundation for such integration.

\paragraph{Evidence composition.}
Composing evidence from multiple sources relates to
Dempster-Shafer theory~\cite{shafer1976mathematical} and
belief revision~\cite{gardenfors1988knowledge}.  The
\JuGeo{} trust lattice provides a simpler, verification-specific
composition mechanism.

\paragraph{Channel architectures.}
The channel concept in \JuGeo{} extends message-passing
architectures~\cite{milner1999communicating} to the evidence
domain with trust annotations.

% ─────────────────────────────────────────────────────────────────────
\section{Conclusion}
\label{sec:conclusion}

We have introduced the Copilot Oracle Channel, a federated evidence
architecture that composes Copilot and solver evidence within the
\JuGeo{} trust framework.  The key results are:

\begin{enumerate}
\item \textbf{No-escalation} (Theorem~\ref{thm:no-escal}):
  federation never produces trust above the minimum ceiling.

\item \textbf{Commutativity} (Theorem~\ref{thm:fed-comm}):
  federation order does not matter.

\item \textbf{Associativity} (Theorem~\ref{thm:fed-assoc}):
  multi-channel federation is well-defined.

\item \textbf{Support route soundness}
  (Theorem~\ref{thm:route-sound}): validated proposals achieve
  solver-level trust.

\item \textbf{Router correctness} (Theorem~\ref{thm:router}):
  only admissible channels are selected.
\end{enumerate}

The oracle pattern demonstrates that LLMs and solvers can be
productively combined within a trust-calibrated framework.

% ─────────────────────────────────────────────────────────────────────
\begin{thebibliography}{25}

\bibitem{young-jugeo}
H.~Young.
\newblock Judgment Geometry: Semantic sites, descent, and the
  Copilot channel.
\newblock JuGeo Paper~0, 2024.

\bibitem{young-channels}
H.~Young.
\newblock The Copilot evidence channel.
\newblock JuGeo Paper~88, 2024.

\bibitem{young-oracle}
H.~Young.
\newblock Copilot oracle federation.
\newblock JuGeo Paper~89, 2024.

\bibitem{young-trust}
H.~Young.
\newblock The trust algebra.
\newblock JuGeo Paper~4, 2024.

\bibitem{young-construction}
H.~Young.
\newblock The Copilot construction participant.
\newblock JuGeo Paper~74, 2024.

\bibitem{chen2021codex}
M.~Chen, J.~Tworek, H.~Jun, et~al.
\newblock Evaluating large language models trained on code.
\newblock \emph{arXiv:2107.03374}, 2021.

\bibitem{first2023baldur}
E.~First, M.~Rabe, T.~Ringer, and Y.~Brun.
\newblock {Baldur}: Whole-proof generation and repair with large
  language models.
\newblock In \emph{Proc.\ ESEC/FSE}, 2023.

\bibitem{yang2024leandojo}
K.~Yang, A.~Swope, A.~Gu, R.~Chaez, and J.~Cai.
\newblock {LeanDojo}: Theorem proving with retrieval-augmented
  language models.
\newblock In \emph{Proc.\ NeurIPS}, 2023.

\bibitem{demoura2008z3}
L.~de~Moura and N.~Bj{\o}rner.
\newblock {Z3}: An efficient {SMT} solver.
\newblock In \emph{Proc.\ TACAS}, 2008.

\bibitem{clarke2003counterexample}
E.~Clarke, O.~Grumberg, S.~Jha, Y.~Lu, and H.~Veith.
\newblock Counterexample-guided abstraction refinement.
\newblock \emph{JACM}, 50(5):752--794, 2003.

\bibitem{hoare1969axiomatic}
C.~A.~R.~Hoare.
\newblock An axiomatic basis for computer programming.
\newblock \emph{CACM}, 12(10):576--580, 1969.

\bibitem{floyd1967assigning}
R.~W.~Floyd.
\newblock Assigning meanings to programs.
\newblock In \emph{Proc.\ Symp.\ Applied Mathematics}, 1967.

\bibitem{shafer1976mathematical}
G.~Shafer.
\newblock \emph{A Mathematical Theory of Evidence}.
\newblock Princeton University Press, 1976.

\bibitem{gardenfors1988knowledge}
P.~G{\"a}rdenfors.
\newblock \emph{Knowledge in Flux}.
\newblock MIT Press, 1988.

\bibitem{milner1999communicating}
R.~Milner.
\newblock \emph{Communicating and Mobile Systems: The
  $\pi$-Calculus}.
\newblock Cambridge University Press, 1999.

\bibitem{xu2008satzilla}
L.~Xu, F.~Hutter, H.~H.~Hoos, and K.~Leyton-Brown.
\newblock {SATzilla}: Portfolio-based algorithm selection for
  {SAT}.
\newblock \emph{JAIR}, 32:565--606, 2008.

\bibitem{wintersteiger2009portfolio}
C.~M.~Wintersteiger, Y.~Hamadi, and L.~de~Moura.
\newblock A concurrent portfolio approach to {SMT} solving.
\newblock In \emph{Proc.\ CAV}, 2009.

\bibitem{de2015lean}
L.~de~Moura, S.~Kong, J.~Avigad, F.~van~Doorn, and
  J.~von~Raumer.
\newblock The {Lean} theorem prover (system description).
\newblock In \emph{Proc.\ CADE}, 2015.

\bibitem{ouyang2022training}
L.~Ouyang, J.~Wu, X.~Jiang, et~al.
\newblock Training language models to follow instructions with
  human feedback.
\newblock In \emph{Proc.\ NeurIPS}, 2022.

\bibitem{dijkstra1975guarded}
E.~W.~Dijkstra.
\newblock Guarded commands, nondeterminacy and formal derivation
  of programs.
\newblock \emph{CACM}, 18(8):453--457, 1975.

\bibitem{leino2010dafny}
K.~R.~M.~Leino.
\newblock {Dafny}: An automatic program verifier for functional
  correctness.
\newblock In \emph{Proc.\ LPAR}, 2010.

\end{thebibliography}

\appendix
\section{Additional Formal Properties}
\label{app:formal}

\subsection{Federation Idempotence}
\label{app:idemp}

\begin{lemma}[Self-federation is identity]
\label{lem:self-fed}
$\OracleComp(s, s) = (w, \tau, \rho, \ell)$ where $w, \tau, \rho, \ell$
are the components of $s$.
\end{lemma}

\begin{proof}
$w \cup w = w$, $\min(\tau, \tau) = \tau$,
$\rho \cap \rho = \rho$, $\max(\ell, \ell) = \ell$.
\end{proof}

\subsection{Trust Composition Under Chains}
\label{app:chains}

\begin{lemma}[Chain trust is the infimum]
\label{lem:chain-trust}
For a federation chain $s_1, s_2, \ldots, s_n$:
$\tau(\OracleComp(s_1, \ldots, s_n)) = \inf(\tau(s_1), \ldots, \tau(s_n))$.
\end{lemma}

\begin{proof}
By induction on $n$.  Base case ($n=1$): trivial.  Inductive step:
$\tau(\OracleComp(s_1, \ldots, s_{n+1}))
= \min(\tau(\OracleComp(s_1, \ldots, s_n)), \tau(s_{n+1}))
= \min(\inf(\tau_1, \ldots, \tau_n), \tau_{n+1})
= \inf(\tau_1, \ldots, \tau_{n+1})$.
\end{proof}

\subsection{Support Route Transitivity}
\label{app:route-trans}

\begin{definition}[Route chain]
\label{def:route-chain}
A \emph{route chain} is a sequence of support routes
$\sigma_1, \sigma_2, \ldots, \sigma_k$ where the target of
$\sigma_i$ matches the source kind of $\sigma_{i+1}$.
\end{definition}

\begin{lemma}[Route chains compose]
\label{lem:route-compose}
If $\sigma_1, \sigma_2$ form a valid route chain, then their
composition $\sigma_2 \circ \sigma_1$ is a valid support route
with target $\sigma_2.t$ and source kind $\sigma_1.k$.
\end{lemma}

\begin{proof}
The composed route inherits the source kind from $\sigma_1$
and the target from $\sigma_2$.  The support region is the
intersection of the two regions.  Validity follows from the
validity of the components and the matching condition.
\end{proof}

\subsection{Router Totality}
\label{app:router}

\begin{lemma}[Router always terminates]
\label{lem:router-term}
For any request $q$, the router terminates in at most $|P|$
steps, where $P$ is the set of registered channels.
\end{lemma}

\begin{proof}
The router iterates over the registered channels (in priority
order), checking jurisdiction for each.  Since $P$ is finite,
the iteration terminates.
\end{proof}

\begin{lemma}[Router preserves preferences]
\label{lem:router-pref}
If the preferred channel $C_p$ is admissible for request $q$,
the router selects $C_p$.
\end{lemma}

\begin{proof}
The router checks the preferred channel first.  If admissible
($J(C_p) \vdash q$), it is selected immediately without
consulting other channels.
\end{proof}

\subsection{No-Escalation Preservation Under Dynamic Federation}
\label{app:dynamic}

\begin{definition}[Dynamic federation]
\label{def:dynamic-fed}
A \emph{dynamic federation} is a sequence of federations
$\mathcal{F}_1, \mathcal{F}_2, \ldots$ where channels may be
added or removed between steps.
\end{definition}

\begin{theorem}[No-escalation is preserved dynamically]
\label{thm:dynamic-noescal}
If each $\mathcal{F}_i$ satisfies the no-escalation invariant,
then the composed result over any prefix of the sequence also
satisfies no-escalation.
\end{theorem}

\begin{proof}
By induction on the sequence length.  The base case is
Theorem~\ref{thm:no-escal}.  For the inductive step, composing
the result of $\mathcal{F}_1, \ldots, \mathcal{F}_n$ (which
satisfies no-escalation by hypothesis) with $\mathcal{F}_{n+1}$
(which satisfies no-escalation by assumption) produces a result
whose trust is the meet of the two, which is at most the
minimum of the two ceilings.
\end{proof}

\subsection{Route Composition}
\label{app:route-comp}

\begin{lemma}[Support route composition is associative]
\label{lem:route-assoc}
For support routes $r_1 : A \to B$, $r_2 : B \to C$, $r_3 : C \to D$:
\[
  (r_3 \circ r_2) \circ r_1 = r_3 \circ (r_2 \circ r_1).
\]
\end{lemma}

\begin{proof}
Route composition is defined as concatenation of trust-annotated
segments.  Concatenation of sequences is associative.  The trust
of the composed route is the infimum of the segment trusts.  The
infimum operation is associative.  Therefore route composition is
associative.
\end{proof}

\begin{lemma}[Empty route is the identity]
\label{lem:route-id}
The empty route $\varepsilon$ satisfies $r \circ \varepsilon = r$
and $\varepsilon \circ r = r$ for all routes~$r$.
\end{lemma}

\begin{proof}
Concatenation with the empty sequence is the identity on sequences.
The infimum of the empty set is $\top = \Tproof$, so composing
with the empty route does not lower the trust.
\end{proof}

\begin{corollary}[Routes form a category]
\label{cor:route-cat}
The set of support routes, with composition as defined above and
the empty route as identity, forms a category.
\end{corollary}

\begin{proof}
Associativity: Lemma~\ref{lem:route-assoc}.
Identity: Lemma~\ref{lem:route-id}.
\end{proof}

\subsection{Federation and Jurisdiction Interaction}
\label{app:fed-juris}

\begin{definition}[Jurisdictional federation]
\label{def:juris-fed}
A federation $\mathcal{F}$ is \emph{jurisdictional} if every
contributing channel operates within its declared jurisdiction:
$\forall C_i \in \mathcal{F}.\; \mathrm{query}(C_i) \subseteq J(C_i)$.
\end{definition}

\begin{lemma}[Jurisdictional federation preserves no-escalation]
\label{lem:juris-noescal}
If $\mathcal{F}$ is jurisdictional, then no-escalation holds.
\end{lemma}

\begin{proof}
Each channel $C_i$ produces evidence within its jurisdiction,
which has trust at most $\tau(C_i)$.  The federation composes
evidence whose trust is $\inf_i \tau(C_i)$, which cannot exceed
any individual $\tau(C_i)$.  Therefore no channel's ceiling is
exceeded.
\end{proof}

\begin{definition}[Oracle scope]
\label{def:oracle-scope}
The \emph{oracle scope} of a Copilot Oracle Channel is the
intersection of the Copilot channel's jurisdiction and the
solver channel's jurisdiction:
$J_{\mathrm{oracle}} = J(\CopChan) \cap J(\SolChan)$.
\end{definition}

\begin{lemma}[Oracle scope is non-empty]
\label{lem:oracle-nonempty}
If both channels have universal jurisdiction (\ie{} they accept
all proposition kinds), then $J_{\mathrm{oracle}} = J(\CopChan) = J(\SolChan)$.
\end{lemma}

\begin{proof}
$J(\CopChan) \cap J(\SolChan) = J(\CopChan)$ when
$J(\SolChan) \supseteq J(\CopChan)$, which holds when
both are universal.
\end{proof}

\begin{remark}[Practical oracle scope]
In practice, the oracle scope may be restricted to
proposition kinds that the solver can decide (e.g., quantifier-free
linear arithmetic), excluding those that require human
judgment.  This restricts the federation to propositions
where solver corroboration is feasible.
\end{remark}

\end{document}
